% =========================================================
% Conclusions and Limitations
% =========================================================

\section{Conclusions and Limitations}

\subsection{Summary of Results}

The present work develops a conditional operator-theoretic perturbative framework associated with a representative bounded-curvature vacuum geometry compatible with the branch structure established in TIG2 \cite{tig2}.

The framework is constructed around the representative lapse function

\begin{equation}
F(r)
=
1
-
\frac{2Mr^2}{r^3+r_c^3},
\end{equation}

which preserves Schwarzschild-compatible asymptotic behavior while introducing a finite representative core scale \cite{wald1984,chandrasekhar1983}.

The horizon sector reduces exactly to the cubic TIG2 branch relation

\begin{equation}
x^3 - x^2 + \beta^3 = 0,
\end{equation}

with critical branch-merger structure

\begin{equation}
x_c = \frac{2}{3},
\qquad
\beta_c^3 = \frac{4}{27}.
\end{equation}

The perturbative dynamics are modeled through a reduced radial operator framework associated with a candidate effective perturbative potential.

The resulting construction yields:
\begin{itemize}[leftmargin=0.6cm]
\item a representative bounded-curvature geometry,
\item a perturbative radial operator structure,
\item a quadratic-form framework,
\item a conditional self-adjointness program,
\item and a conditional spectral admissibility hierarchy.
\end{itemize}

% =========================================================

\subsection{Interpretation of the Present Framework}

The present work should be interpreted conservatively as a structured perturbative admissibility program within mathematical physics.

The framework does not attempt to establish a complete gravitational theory.

Rather, the objective is to define a mathematically coherent vacuum-sector structure suitable for further operator-theoretic and geometric analysis.

In particular, the present analysis develops:
\begin{itemize}[leftmargin=0.6cm]
\item perturbative consistency structure,
\item admissibility conditions,
\item and conditional bounded linear evolution.
\end{itemize}

The framework therefore represents a mathematically controlled perturbative sector rather than a final gravitational completion.

% =========================================================

\subsection{Conditional Nature of the Results}

Several central components of the framework remain conditional.

In particular:
\begin{itemize}[leftmargin=0.6cm]
\item the exact derivation of the effective perturbative potential,
\item endpoint classification,
\item quadratic-form positivity,
\item and full self-adjoint realization
\end{itemize}

remain open operator-theoretic problems \cite{friedrichs1934,reedsimon2,reedsimon4,kato1995,weidmann1987}.

Accordingly, the spectral admissibility condition

\begin{equation}
\sigma(L_{\mathrm{TIG}}^{F})
\subset
[0,\infty)
\end{equation}

is presently interpreted as a conditional admissibility requirement rather than a rigorously established theorem.

% =========================================================

\subsection{Non-Claims}

The present work does not establish:
\begin{itemize}[leftmargin=0.6cm]
\item nonlinear stability,
\item complete spectral classification,
\item full covariant field equations,
\item quantum-gravity completion,
\item observational verification,
\item or complete singularity-resolution theorems.
\end{itemize}

The framework also does not claim to replace General Relativity.

Instead, the representative geometry is constructed specifically to preserve Schwarzschild-compatible asymptotics within the low-curvature regime \cite{wald1984,hawkingellis1973}.

% =========================================================

\subsection{Near-Critical Sector}

Near the critical branch-merger regime

\begin{equation}
\beta \approx \beta_c,
\end{equation}

the representative geometry may develop modified perturbative spectral structure associated with elongated throat behavior and metastable trapping sectors.

Potential effects include:
\begin{itemize}[leftmargin=0.6cm]
\item resonance accumulation,
\item slow decay,
\item spectral compression,
\item and long-lived perturbative sectors.
\end{itemize}

The present work does not establish rigorous resonance theorems in this regime.

Such structures remain conjectural and require separate asymptotic spectral analysis.

% =========================================================

\subsection{Scientific Position}

The strongest currently justified interpretation of the present framework is the following:

\begin{quote}
TIG3 provides a coherent conditional operator-theoretic perturbative admissibility framework for a representative bounded-curvature vacuum sector compatible with the TIG2 cubic branch structure.
\end{quote}

This statement defines the intended scientific scope of the present work.

% =========================================================

\subsection{Future Directions}

Several mathematically significant problems remain open.

These include:
\begin{itemize}[leftmargin=0.6cm]
\item rigorous derivation of the effective perturbative potential,
\item endpoint classification,
\item quadratic-form coercivity,
\item spectral nonnegativity,
\item resonance analysis in the near-critical regime,
\item and possible nonlinear extensions of the perturbative framework.
\end{itemize}

Future work may also investigate whether the present perturbative admissibility structure can be embedded into a more complete covariant geometric formulation.

Such developments remain outside the scope of the present paper.

% =========================================================

\subsection{Final Remarks}

The present work establishes a mathematically conservative perturbative admissibility framework associated with a representative bounded-curvature vacuum geometry compatible with the TIG2 branch structure.

The framework is intended as a structurally controlled mathematical physics program emphasizing:
\begin{itemize}[leftmargin=0.6cm]
\item explicit assumptions,
\item theorem transparency,
\item operator-theoretic consistency,
\item and publication-stable scientific interpretation.
\end{itemize}

The present analysis therefore represents a structured continuation program in perturbative geometric operator theory rather than a final theory of gravitation.
